% ============ 第1章 参考文献(原书 pp.122-123) ============
\section*{参考文献}\addcontentsline{toc}{section}{参考文献}
\begin{thebibliography}{99}
% 日文条目使用日文字体(Fandol 字体无日文假名/和制汉字字形;用 TeX Live 自带的 Harano Aji Mincho)
\newfontfamily\jpref{Harano Aji Mincho}
\item[] 这里,列举本章参考的文献以及进一步深化本章内容的文献.
\item[] 单行本有:
\bibitem[1]{ch1:1} {\jpref 猪狩惺, 実解析入門, 岩波書店, 東京, 1996.}
\bibitem[2]{ch1:2} {\jpref 伊藤清三, ルベーグ積分入門, 裳華房, 東京, 1963.}
\bibitem[3]{ch1:3} R. R. Coifman and Y. Meyer, Au-delà des opérateurs pseudodifférentiels, Astérisque \#57, Société Math. France, 1978.
\bibitem[4]{ch1:4} J. García-Cuerva and J. L. Rubio de Francia, Weighted Norm Inequalities and Related Topics, North-Holland Math. Studies 116, Amsterdam, 1985.
\bibitem[5]{ch1:5} L. Grafakos, Classical and Modern Fourier Analysis, Pearson Education, NJ, 2004.
\bibitem[6]{ch1:6} J.-L. Journé, Calderón--Zygmund Operators, Pseudo-Differential Operators and the Cauchy Integral of Calderón, Lecture Notes in Math. \#994, Springer-Verlag, Berlin, 1983.
\bibitem[7]{ch1:7} E. M. Stein, Singular Integrals and Differentiability Properties of Functions, Princeton University Press, Princeton, NJ, 1970.
\bibitem[8]{ch1:8} E. M. Stein and G. Weiss, Introduction to Fourier Analysis on Euclidean Spaces, Princeton University Press, Princeton, NJ, 1971.
\bibitem[9]{ch1:9} E. M. Stein, Harmonic Analysis: Real Variable Methods, Orthogonality, and Oscillatory Integrals, Princeton University Press, Princeton, NJ, 1993.
\bibitem[10]{ch1:10} A. Torchinsky, Real-Variable Methods in Harmonic Analysis, Academic Press, New York, 1986.
\bibitem[11]{ch1:11} A. Uchiyama, Hardy Spaces on the Euclidean Space, Springer-Verlag, Tokyo, 2001.
\bibitem[12]{ch1:12} A. Zygmund, Trigonometric Series, I and II, 2nd Edition, Cambridge University Press, Cambridge, UK, 1959.
\item[] 第 1.1 节用文献 \cite{ch1:1}, \cite{ch1:2}, \cite{ch1:3}.第 1.2,1.3 节用文献 \cite{ch1:8}.第 1.4～1.8 节用文献 \cite{ch1:4}, \cite{ch1:5}, \cite{ch1:6}, \cite{ch1:7}, \cite{ch1:10}.第 1.8 节中 Marcinkiewicz 插值定理的证明参考了文献 \cite{ch1:12}.第 1.9 节用文献 \cite{ch1:3},第 1.10 节用文献 \cite{ch1:5},第 1.11 节用文献 \cite{ch1:4}, \cite{ch1:5}.关于 $H^{1}(\mathbb{R})$,$\mathrm{BMO}(\mathbb{R})$,与论文 C. Fefferman and E. M. Stein, $H^{p}$ spaces of several variables, Acta Math. 129 (1972), 137-193. 一起参考了文献 \cite{ch1:11}.
\item[] 想在实分析领域进一步扩展知识的人,不妨翻阅文献 \cite{ch1:5}, \cite{ch1:7}, \cite{ch1:9}.关于 Hardy 空间,BMO 空间,文献 \cite{ch1:11} 很详细.
\end{thebibliography}